%\documentclass[11pt,a4paper,twoside]{article}

\usepackage[utf8]{inputenc}

\usepackage[T1,T2A]{fontenc}

\usepackage[english.ancient,russian.ancient]{babel}

\usepackage{graphicx}

\usepackage{array}

\usepackage[doi]{samgtu-bib}

\usepackage{samgtu}

\usepackage{samgtu-name}

\usepackage{mathtools}

\mathtoolsset{showonlyrefs}

\usepackage{desclist}

\usepackage{euscript}

\usepackage{mathrsfs} 

\usepackage{multicol}

\usepackage{mathrsfs}

\usepackage{cite}

\usepackage{cmap}

\usepackage{qrcode}

\allowdisplaybreaks[4]

\usepackage[nodayofweek]{datetime}

\usepackage[thinlines]{easybmat}



\renewcommand{\JournalYear}{2018}

\renewcommand{\JournalVolume}{22}

\renewcommand{\JournalNumber}{4}



\newdate{JournalDateSign}{29}{12}{2018}% Дата подписания Вестника в печать

\renewcommand{\JournalDateSign}{\displaydate{JournalDateSign}}

\newdate{JournalDatePrint}{16}{01}{2019}% Дата выхода Вестника из печати

\renewcommand{\JournalDatePrint}{\displaydate{JournalDatePrint}}

%\renewcommand{\JournalUPL}{16.56} % Усл. печ. л.

%\renewcommand{\JournalUIL}{16.53} % Уч.-изд. л.

\renewcommand{\JournalUPL}{15.24} % Усл. печ. л.

\renewcommand{\JournalUIL}{15.21} % Уч.-изд. л.

\renewcommand{\JournalRegNumber}{220/18} % Регистрационный номер

\renewcommand{\JournalOrderNumber}{2} % Номер заказа



%\renewcommand{\JournalRMonth}{Март} % Месяц выхода Вестника

%\renewcommand{\JournalEMonth}{March} % Месяц выхода Вестника

%\renewcommand{\JournalRMonth}{Июнь} % Месяц выхода Вестника

%\renewcommand{\JournalEMonth}{June} % Месяц выхода Вестника

%\renewcommand{\JournalRMonth}{Сентябрь} % Месяц выхода Вестника

%\renewcommand{\JournalEMonth}{September} % Месяц выхода Вестника

\renewcommand{\JournalRMonth}{Декабрь} % Месяц выхода Вестника

\renewcommand{\JournalEMonth}{December} % Месяц выхода Вестника



\newcommand{\totop}[1]{\raisebox{6pt}[1pt][2pt]{#1}}

\newcommand{\tottop}[1]{\raisebox{12pt}[1pt][2pt]{#1}} 

\newcommand{\totttop}[1]{\raisebox{18pt}[1pt][2pt]{#1}} 

\newcommand{\tottttop}[1]{\raisebox{24pt}[1pt][2pt]{#1}} 

\newcommand{\totttttop}[1]{\raisebox{30pt}[1pt][2pt]{#1}} 

\newcommand{\tobot}[1]{\raisebox{-6pt}[1pt][2pt]{#1}} 

\newcommand{\tobbot}[1]{\raisebox{-12pt}[1pt][2pt]{#1}} 

\newcommand{\tobbbot}[1]{\raisebox{-18pt}[1pt][2pt]{#1}}



\graphicspath{{./00PIC/}}

%\usepackage{samgtu-name}

%\usepackage{cmap}

%

%%

\vsgtu{1585} %registration number



\FirstPage{225}

\LastPage{235}

%%

\ResearchArticle

%%

%\pdfcrop

%\draft

%%

%

%\begin{document}

%

%

\hfill \qrcode[hyperlink,height=.8in]{http://doi.org/10.14498/vsgtu1585}

\vspace{-.8in}





\udc{517.955.2:517.956}

\msc{35M12, 35A02}



\rutitle{Критерий однозначной разрешимости спектральной задачи Дирихле для одного класса многомерных гиперболо-параболических   уравнений}

\rutitleshort{Критерий однозначной разрешимости спектральной задачи Дирихле\dots}

%\rutitlehack{Задача Коши для системы дифференциальных уравнений гиперболического типа порядка $\boldsymbol n$ с~некратными характеристиками}



\entitle{A criterion for the unique solvability of the spectral Dirichlet problem for a class of multidimensional hyperbolic-parabolic equations}

\entitleshort{A criterion for the unique solvability of the spectral Dirichlet problem\dots}

%\entitlehack{The Cauchy problem for a system of the hyperbolic

%differential equations of the $\boldsymbol n$-th order with~the~nonmultiple~characteristics}



\ruauthor{С.~А.~Алдашев}{А\,л\,д\,а\,ш\,е\,в~С.~А.}

\enauthor{S.~A.~Aldashev}{A\,l\,d\,a\,s\,h\,e\,v~S.~A.}



\ruaffil{\ruaffid{2823}{Казахский национальный педагогический университет им. Абая,\\

Казахстан, 480100, Алматы, ул. Толе Би, 86}.}

\enaffil{\enaffid{2823}{Kazakh National Pedagogical University named after Abay,\\

86, Tole-bi st.,  Almaty, 480100, Kazakhstan}.}









\ruauthorsdetails{1}{% Количество авторов

\fullauthor{\rupid{43939}{Серик Аймурзаевич Алдашев}}

\CAuthor % Автор-корреспондент

\orcid{0000-0002-8223-6900} % ORCID ID автора 

\regalia{доктор физико-математических наук, профессор} 

\position{заведующий кафедрой} 

\dept{каф. фундаментальной и прикладной математики} 

\email{aldash51@mail.ru}

}



\enauthorsdetails{1}{

\fullauthor{\enpid{43939}{Serik A. Aldashev}}

\CAuthor % Автор-корреспондент

\orcid{0000-0002-8223-6900} % ORCID ID автора 

\regalia{Dr. Phys. \& Math. Sci., Professor} 

\position{Head of Dept.} 

\dept{Dept. of Fundamental and Applied Mathematics}

\email{aldash51@mail.ru}

}



\rukeywords{многомерное гиперболо-параболическое уравнение, спектральная задача Дирихле, многомерная цилиндрическая область}



\enkeywords{multidimensional hyperbolic-parabolic equation, Dirichlet spectral problem, multidimensional cylindrical domain}





\ruabstract{В цилиндрической области евклидова пространства для одного класса многомерного

гиперболо-параболических уравнений рассматривается

спектральная задача Дирихле с однородными краевыми условиями. Решение ищется в виде разложения по многомерным сферическим функциям. Доказаны теоремы существования и единственности решения. Получены условия однозначной разрешимости поставленной задачи, которые существенно зависят от высоты цилиндра.}



\enabstract{In the cylindrical domain of Euclidean space for one class of multidimensional hyperbolic parabolic equations the spectral Dirichlet problem with homogeneous boundary conditions is considered. The solution is sought in the form of an expansion in multidimensional spherical functions. The existence and uniqueness theorems of the solution are proved. Conditions for the unique solvability of the problem are obtained, which essentially depend on the height of the cylinder.}



\newdate{aldasha}{05}{12}{2017}

\date{\displaydate{aldasha}}

\newdate{aldashb}{08}{04}{2018}

\revisiondate{\displaydate{aldashb}}

\newdate{aldashc}{11}{06}{2018}

\accepteddate{\displaydate{aldashc}}



\newdate{aldashe}{01}{07}{2018}

\renewcommand{\JournalDataOnline}{\displaydate{aldashe}}



\vspace{5mm}



%\ForPeerReview

\makerutitle



\renewcommand{\baselinestretch}{0.93}



\Section[n]{Введение} 

Двумерные спектральные задачи для гиперболических уравнений интенсивно изучаются, а их многомерные аналоги, насколько известно автору, исследованы мало. Это связано с тем, что в случае трех и~более независимых переменных возникают трудности принципиального характера, так как весьма привлекательный и удобный метод сингулярных интегральных уравнений, применяемый для двумерных задач, здесь не может быть использован из-за отсутствия сколько-нибудь полной теории многомерных сингулярных интегральных уравнений. Теория многомерных сферических функций, напротив, достаточно полно изучена. Эти функции имеют важные приложения в математической физике, в~теоретической физике и~в~теории многомерных сингулярных интегральных уравнений. Автор предлагает при решении спектральных задач Дирихле для многомерных гиперболо"=параболических уравнений использовать разложения по сферическим функциям.



\bigskip



\textbf{1. Постановка задачи и основной результат.} 

Краевые задачи для гиперболо"=параболических уравнений на

плоскости хорошо изучены~\cite{aldash:1}. 

По мнению автора, их многомерные аналоги  \cite{aldash:2, aldash:3, aldash:4} исследованы недостаточно полно.



Пусть $\Omega_{\alpha\beta} $  "--- цилиндрическая область евклидова

пространства $E_{m+1}$ точек $(x_{1}, \dots, x_{m} ,t),$ ограниченная

цилиндром $\Gamma = \{(x,t):  |x|=1 \}$, плоскостями

$t=\alpha>0$ и $t=\beta<0,$ где ${\left| {x}\right|} $ "---  длина

вектора $x = (x_{1},\dots ,x_{m}).$



Обозначим через  $\Omega_\alpha$ и $\Omega_\beta$  части области

$\Omega_{\alpha \beta},$ а через $\Gamma_\alpha$,  $\Gamma_\beta$ "--- части поверхности $\Gamma,$  лежащие в

полупространствах $t>0$ и  $t<0$; $\sigma_\alpha$ "--- верхнее, а

$\sigma_\beta$ "--- нижнее основание области $\Omega_{\alpha \beta}.$



Пусть далее $S$ "--- общая часть границ областей $\Omega_\alpha$,

$\Omega_\beta$ представляющая множество $\{t=0, \, 0<|x|<1\}$  в~$E_m.$



В области $\Omega_{\alpha\beta}$ рассмотрим многомерные смешанные гиперболо"=параболические уравнения со спектральным действительным параметром $\gamma$:

\begin{equation}

\gamma u =

\left\{\begin{array}{ll}\Delta_x u- u_{tt} + \sum\limits_{i=1}^m a_i(x,t)u_{x_i}+ b(x,t)u_t + c(x,t)u,& t>0,\\

\Delta_x u - u_{t}+ \sum\limits_{i=1}^m d_i(x,t)u_{x_i}+

e(x,t)u, & t<0, \end{array} \right.

\label{ald:eq1}

\end{equation}

где $\Delta_{x}$ "--- оператор Лапласа по переменным $x_{1} ,\dots, x_{m}$, $m \ge 2.$



В качестве многомерной спектральной задачи Дирихле рассмотрим

следующую задачу.



\bigskip



\hypertarget{ald:taskD}{}



{\small\sc Задача D.} \em Найти решение уравнения \eqref{ald:eq1}  в области

$\Omega_{\alpha\beta}$ при $t\neq0$ из класса

$C(\overline{\Omega}_{\alpha\beta}) \cap C^2 (\Omega_\alpha\cup

\Omega_\beta),$ удовлетворяющее краевым условиям

\begin{gather}

  u\big|_{\sigma_\alpha} = 0,\quad

  u\big|_{\Gamma_\alpha}= 0,

\label{ald:eq2}

\\

\label{ald:eq3}

u\big|_{\Gamma_\beta}= 0,\quad  u\big|_{\sigma_\beta} = 0. 

\end{gather}



\em



\bigskip



Для  удобства перейдем от декартовых координат

$x_1,\dots,x_m,t$ к сферичес\-ким $r, \theta_1,\dots,\theta_{m-1},t, $ 

$r\ge 0,$  $ 0\le \theta_1 < 2\pi$, $0 \le \theta_i \leq \pi$, 

$i= 2,3,\dots,m-1$, $\theta=(\theta_1,\dots,\theta_{m-1}).$



Пусть $\left\{ {Y_{n,m}^{k}(\theta )}\right\} $ "--- система линейно

независимых сферических функций порядка $n$, $1 \le k \le k_{n}$,

$(m - 2)!n!k_{n} = (n + m - 3)!(2n + m - 2)$,

$W_{2}^{l}(S)$, $l = 0,1,\dots $ "--- пространства Соболева.



Через 

$\widetilde{d}_{in}^k(r,t)$, 

$d_{in}^k(r,t)$,

$\widetilde{e}_n^k(r,t)$,

$\widetilde{d}_n^k(r,t)$, 

$\rho_n^k$  обозначим коэффициенты разложения рядов по сферическим функциям

$Y_{n,m}^{k}(\theta )$ соответственно функций

$d_i(r,\theta,t)\rho$,

$d_i\frac{x_i}{r}\rho$, 

$e(r,\theta,t)\rho$, 

$d(r,\theta,t)\rho$, 

$\rho(\theta)$, $i=1,2, \dots , m,$

причем $\rho(\theta)\in C^\infty (H),$ $H$ "--- единичная сфера в~$E_m.$



Пусть $a_i(x,t)$, $b(x,t)$, $c(x,t)\in W_2^l(\Omega_\alpha)\subset C

(\bar{\Omega}_\alpha)$, $d_i(x,t)$, $e(x,t)\in W_2^l(\Omega_\beta)$, $i = 1,2,\dots ,m$, 

$ l\geq m+1$, $e(x,t)-\gamma\leq 0$, $\forall(x,t)\in \Omega_\beta.$









\pagebreak



\hypertarget{ald:th1}{}



{\small\sc Теорема 1.}  \em

Справедливы следующие утверждения\em. \em

\begin{itemize}

\item[\rm 1.] Если $\gamma \leq -\mu_{s,n}^2,$ то задача~\hyperlink{ald:taskD}{{\small \sc D}}

имеет только тривиальное $($нулевое$)$ решение\em. \em 

\item[\rm 2.] При $\gamma > -\mu_{s,n}^2$  задача~\hyperlink{ald:taskD}{{\small \sc D}} имеет только тривиальное решение тогда и~только тогда\em, \em когда

 $$

\sin \alpha \sqrt{\gamma+\mu_{s,n}^2}\neq 0,\quad s=1, 2,\dots,

 $$

где  $\mu_{s,n}  $ "--- нули функций Бесселя первого рода $J_{n+{(m-2)}/{2}}(z).$ \rm



\end{itemize}



\bigskip \rm



Отметим, что эта теорема при $a_i(x,t)= b(x,t)= c(x,t) = d_i(x,t)\hm= e(x,t)

=0$, $i=1,2, \dots,n,$  сформулирована в~\cite{aldash:5}.



\bigskip



\hypertarget{ald:p2}{}



\textbf{2. Разрешимость задачи~\hyperlink{ald:taskD}{D}.} 

В~сферических координатах уравнение \eqref{ald:eq1} в~области $\Omega_\beta$ имеет вид

\begin{equation}

L_1 u\equiv u_{rr}+\dfrac{m-1}{r}u_r - \dfrac{1}{r^2}\delta u -

u_t + \sum\limits_{i=1}^m d_i(r,\theta,t)u_{x_i}+ e(r,\theta,t)u =

\gamma u, 

\label{ald:eq4}

\end{equation}

где

$$

\delta \equiv -\sum\limits_{j=1}^{m-1} \dfrac{1}{g_j \sin^{m-j-1}\theta_j}\dfrac{\partial}{\partial

\theta_j}

\Bigl(\sin^{m-j-1}\theta_j\dfrac{\partial}{\partial \theta_j}\Bigr),

$$

$$

g_1=1,\quad 

g_j=(\sin\theta_1 \cdots \sin\theta_{j-1})^2,\quad  j>1.

$$



Известно~\cite{aldash:6}, что спектр оператора $\delta$ состоит из

собственных чисел $\lambda_n \hm = n (n+m -2)$, $n=0,1,\dots,$

каждому из которых соответствует $k_n$ ортонормированных

собственных функций $Y_{n,m}^k (\theta).$



Искомое решение задачи \eqref{ald:eq1}, \eqref{ald:eq3} в области $\Omega_\beta$ будем

искать в виде ряда

\begin{equation}

u(r,\theta ,t) = {\sum\limits_{n = 0}^{\infty} {{\sum\limits_{k =

1}^{k_{n} } {\bar {u}_{n}^{k} (r,t)Y_{n,m}^{k} (\theta )}}} } ,

\label{ald:eq5}

\end{equation}

где $\bar {u}_{n}^{k} (r,t)$ "--- функции, подлежащие определению.



Подставив \eqref{ald:eq5}  в \eqref{ald:eq4}, умножив полученное выражение на $\rho(\theta)\neq 0$ и проинтегрировав по единичной сфере $H$ для $\overline{u}_n^k$, получим~\cite{aldash:4}

\begin{multline}

\label{ald:eq6}

\rho_0^1\bar{u}_{0rr}^{1} - \rho_0^1 \bar {u}_{0t}^{1} +

\Bigl(\dfrac{m-1}{r} \rho_0^1 + \sum\limits_{i=1}^m

d_{i0}^1\Bigr)\bar{u}_{0r}^1 + \tilde{e}_0^1 \bar{u}_{0}^1 -

\gamma \rho_0^1 \bar{u}_0^1 + 

\\

+

\sum\limits_{n=1}^\infty

\sum\limits_{k=1}^{k_n} 

\Bigl\{ 

\rho_n^k \bar{u}_{nrr}^{k} - \rho_n^k \bar {u}_{nt}^{k}+ 

\Bigl(\dfrac{m-1}{r}\rho_n^k + \sum\limits_{i=1}^m

d_{in}^k\Bigr)\bar{u}_{nr}^k +

\\ 

+ \Bigl[\tilde{e}_n^k - \lambda_n \dfrac{\rho_n^k}{r^2} +\sum\limits_{i=1}^m(\tilde{d}_{in-1}^k - n

d_{n}^k)\Bigr]\bar{u}_n^k - \gamma \rho_n^k \bar{u}_n^k\Bigr\} =0.

\end{multline}



Теперь рассмотрим бесконечную систему дифференциальных уравнений:

\begin{equation}

\rho_0^1\bar{u}_{0rr}^{1} - \rho_0^1 \bar {u}_{0t}^{1} +

\dfrac{m-1}{r}\rho_0^1\bar{u}_{0r}^{1}= \gamma \rho_0^1

\bar{u}_0^1,

\label{ald:eq7}

\end{equation}

\begin{multline}

\rho_1^k \bar{u}_{1rr}^{k} - \rho_1^k \bar {u}_{1t}^{k} +

\dfrac{m-1}{r} \rho_1^k \bar{u}_{1r}^{k}-

\dfrac{\lambda_1}{r^2}\rho_1^k \bar{u}_1^k =

\\ 

= \gamma \rho_0^1 \bar{u}_1^k -\dfrac{1}{k_1}

\Bigl(\sum\limits_{i=1}^m d_{i0}^1\bar{u}_{0r}^{1} +

\tilde{e}_0^1\bar{u}_{0}^1\Bigr), \quad n=1,\,\, k=1,2,\dots,k_1,

\label{ald:eq8}

\end{multline}

\begin{multline}

\rho_n^k \bar{u}_{nrr}^{k} - \rho_n^k \bar {u}_{nt}^{k} +

\dfrac{m-1}{r} \rho_n^k\bar{u}_{nr}^{k}- \dfrac{\lambda_n}{r^2}

\rho_n^k \bar{u}_n^k =

\\ 

= \gamma \rho_n^k \bar{u}_n^k -\dfrac{1}{k_n}\sum\limits_{k=1}^{k_{n-1}} 

\biggl\{\sum\limits_{i=1}^m d_{in-1}^k \bar{u}_{n-1r}^{k} +

\Bigl[\tilde{e}_{n-1}^k + \sum\limits_{i=1}^m(\tilde{d}_{in-2}^k -

(n-1)d_{n-1}^k)\Bigr]\bar{u}_{n}^k \biggr\},\\ k =1,2,\dots k_n,\,\,n = 2, 3, \dots.

\label{ald:eq9}

\end{multline}



Суммируя уравнения \eqref{ald:eq8} от 1 до $k_1,$ а уравнения \eqref{ald:eq9} "--- от 1 до $k_n,$ а затем складывая полученные выражения вместе с \eqref{ald:eq7}, приходим к~уравнению \eqref{ald:eq6}.



Отсюда следует, что если система функций $\left\{ \bar {u}_{n}^{k}\right\},$

$k=1,2, \dots, k_{n}$, $n \hm = 0,1,\dots,$ "--- решение системы уравнений \eqref{ald:eq7}--\eqref{ald:eq9},

то она является и решением уравнения~\eqref{ald:eq6}.



Нетрудно заметить, что каждое уравнение системы \eqref{ald:eq7}--\eqref{ald:eq9} можно представить в виде

\begin{equation}

\bar {u}_{nrr}^{k} -\bar{u}_{nt}^k + {\dfrac{m - 1}{{r}}}\bar

{u}_{nr}^{k} - {\dfrac{{\lambda_{n}}}{{r^{2}}}}\bar{u}_{n}^{k} -

\gamma \bar{u}_n^k = \bar{f}_{n}^{k}(r,t), 

\label{ald:eq10}

\end{equation}

где $\bar{f}_{n}^{k} (r,t)$ определяются из предыдущих уравнений

этой системы, причем $\bar{f}_{0}^{1}(r,t) \equiv 0.$

Далее из краевого условия \eqref{ald:eq3} в~силу \eqref{ald:eq5} будем иметь 

\begin{equation}

\bar{u}_n^k(r,\beta)= 0,\quad 

\bar{u}_{n}^k(1,t)= 0,\quad  k =1,2, \dots k_{n},\,\, n = 0, 1,\dots.

\label{ald:eq11}

\end{equation}



Произведя в \eqref{ald:eq10}, \eqref{ald:eq11} замену $\bar{u}_{n}^{k} (r,t)= r^{\frac{1-m}{2}}u_{n}^{k} (r,t),$ получим задачу

\begin{gather}

L u_n^k \equiv \bar u_{nrr}^{k} - u_{nt}^k - {\dfrac{\bar \lambda{

_{n}} }{{r^{2}}}}u_{n}^{k} - \gamma u_n^k = f_{n}^{k} (r,t),

\label{ald:eq12}

\\

\label{ald:eq13}

u_n^k(r,\beta)= 0,\quad 

u_{n}^k(1,t)=0,\quad 

 k = 1,2, \dots, k_{n},\, \, n = 0,1,\dots,

\end{gather}

где

$$

\bar{\lambda}_{n} =\dfrac{(m - 1)(3 - m) }{4}- \lambda_{n},

\quad

f_n^k (r,t)= r^{\frac{m-1}{2}}\bar{f}_n^k (r,t).

$$







Решение задачи \eqref{ald:eq12}, \eqref{ald:eq13}  будем искать в виде

\begin{equation}

u_n^k (r,t)=\sum\limits_{s=1}^{\infty} R_s(r) T_s(t), 

\label{ald:eq14}

\end{equation}

при этом пусть

\begin{equation}

f_n^k(r,t)=\sum\limits_{s=1}^{\infty} a_{s,n}^k(t) R_s(r).

\label{ald:eq15}

\end{equation}



Подставляя \eqref{ald:eq14} в \eqref{ald:eq12}, \eqref{ald:eq13} и  учитывая \eqref{ald:eq15}, получим

\begin{gather}

R_{srr}+ \dfrac{\bar{\lambda}_n}{r^2}R_s + (\mu-\gamma) R_{s}=

0,\quad  0<r<1, 

\label{ald:eq16}

\\

R_{s}(1)=0,\quad |R_s(0)|< \infty, 

\label{ald:eq17}

\\

T_{st} + \mu T_s(t) = -a_{s,n}^k(t),\quad \beta<t<0,

\label{ald:eq18}

\\

T_s(\beta) = 0,

\label{ald:eq19}

\end{gather}

где $\mu = \gamma+\mu_{s,n}^2.$



Ограниченным решением задачи \eqref{ald:eq16}, \eqref{ald:eq17} является функция~\cite[c.~401, формула (1{\sl a}\/)]{aldash:7}

\begin{equation}

R_s(r) = \sqrt{r}J_\nu(\mu_{s,n} r),\quad \nu = n+\frac{m-2}{2}.

\label{ald:eq20}

\end{equation}





Решением задачи \eqref{ald:eq18}, \eqref{ald:eq19} является функция~\cite[c.~35, п.~4.3]{aldash:7}

\begin{equation}

T_{s,n}(t) =  \bigl(\exp (-\gamma-\mu_{s,n}^2)t\bigr) 

\int _{t}^{\beta} a_{s,n}^k(\xi) \bigl(\exp (\gamma + \mu_{s,n}^2)\xi \bigr) d \xi.

\label{ald:eq21}

\end{equation}



Подставляя \eqref{ald:eq20} в \eqref{ald:eq15}, будем иметь

\begin{equation}

r^{-\frac{1}{2}} f_n^k(r,t)=\sum\limits_{s=1}^{\infty} a_{s,n}^k(t)

J_{\nu}(\mu_{s,n} r),\quad  0<r<1.

\label{ald:eq22}

\end{equation}



Ряд \eqref{ald:eq22} "--- разложение в ряд Фурье"--~Бесселя \cite{aldash:8}, если

\begin{equation}

a_{s,n}^k(t)=

\dfrac{1}{[J_{\nu+1}(\mu_{s,n})]^{2}}\int _{0}^{1}\sqrt{\xi}f_n^k(\xi,t)J_{\nu}(\mu_{s,n}

\xi)d\xi,

\label{ald:eq23}

\end{equation}

где

$\mu_{s,n}$, $ s= 1,2,\dots$ "--- положительные нули функций Бесселя $J_\nu(z)$, расположенные в порядке возрастания их величины.



Из \eqref{ald:eq20}, \eqref{ald:eq21} получим решение задачи \eqref{ald:eq12}, \eqref{ald:eq13}:

\begin{equation}

u_{n}^k(r,t)=\sum\limits_{s=1}^{\infty} \sqrt{r} T_{s,n}(r) J_{\nu}(\mu_{s,n} r),

\label{ald:eq24}

\end{equation}

где $a_{s,n}(t)$ находятся из \eqref{ald:eq23}. 

Следовательно, сначала решив задачу \eqref{ald:eq7}, \eqref{ald:eq11} $(n=0),$ а затем задачу

\eqref{ald:eq8}, \eqref{ald:eq11} $(n=1)$ и~т.~д., найдем последовательно все $u_n^k (r,t)$

из \eqref{ald:eq24}, $k =1,2,\dots, k_{n}$, $n = 0, 1,\dots.$



Итак, в области $\Omega_\beta$ имеет место равенство

\begin{equation}

\int_{H}\rho(\theta)(L_1-\gamma)udH=0. 

\label{ald:eq25}

\end{equation}



Пусть $f(r,\theta,t)=R(r)\rho(\theta)T(t),$ причем $R(r)\in V_0$,

$V_0 $ плотно в $L_2((0,1))$; 

$\rho(\theta)\in C^{\infty}(H)$  плотно в $L_2(H),$ а 

$T(t)\in V_1$, $V_1$ плотно в~$L_2((\beta,0)).$ 

Тогда $f(r,\theta, t)\in V$, $V=V_0\otimes H \otimes V_1$ плотно в $L_2(\Omega_\beta)$  \cite{aldash:9}.



Отсюда и из \eqref{ald:eq25} следует

$$

\int_{\Omega_\beta}f(r,\theta,t)(L_1-\gamma)ud \Omega_\beta =0

$$

и

$$

L_1 u= \gamma u \quad  \forall(r,\theta,t)\in \Omega_\beta.

$$



Таким образом, решением задачи \eqref{ald:eq1}, \eqref{ald:eq3} в области $\Omega_\beta$

является функция

\begin{equation}

u(r,\theta ,t) = \sum\limits_{n = 0}^{\infty} \sum\limits_{k=1}^{k_{n}} \sum\limits_{n = 0}^{\infty}

r^{\frac{(2-m)}{2}}T_{s,n}(t) J_{n+\frac{(m-2)}{2}}(\mu_{s,n}r)Y_{n,m}^{k} (\theta ),

\label{ald:eq26}

\end{equation}

где $T_{s,n}(t)$ определяются из \eqref{ald:eq21}.



Из \eqref{ald:eq23}, \eqref{ald:eq21}, \eqref{ald:eq24} следует, что $a_{s,n}(t)= T_{s,n}(t)=0$ и

$u_n^k(r,t)=0$, $s\hm =1, 2,\dots$, $k =1,2, \dots, k_{n}$, $n = 0,1, \dots.$



В свою очередь, из \eqref{ald:eq26} получим, что $u(r,\theta,t)=0$ в~$\Omega_\beta.$



Отсюда при  $t\to -0$ получим

\begin{equation}

 u(r,\theta,0) =0.

 \label{ald:eq27}

\end{equation}



Таким образом, учитывая краевые условия \eqref{ald:eq2}, \eqref{ald:eq27}, в

области $\Omega_\alpha$ получим  спектральную задачу Дирихле для уравнения

\begin{equation}

\Delta_x u - u_{tt} + \sum\limits_{i=1}^m a_i(x,t)u_{x_i}+

b(x,t)u_t + c(x,t)u,= \gamma u 

 \label{ald:eq28}

\end{equation}

c условиями

\begin{equation}

u\big|_{S\cup \Gamma_\alpha \cup \sigma_\alpha} = 0.

 \label{ald:eq29}

\end{equation}



В работе  \cite{aldash:10} доказана  справедливость теоремы~\hyperlink{ald:th1}{1} для задачи \eqref{ald:eq28}, \eqref{ald:eq29}.

Следовательно, разрешимость задачи~\hyperlink{ald:taskD}{\small\sc D} установлена.



\bigskip



\textbf{3. Единственность решения задачи~\hyperlink{ald:taskD}{D}.} 

Сначала рассмотрим задачу \eqref{ald:eq1}, \eqref{ald:eq3} в области $\Omega_\beta$ и докажем  единственность ее решения. 

Для этого сначала построим решение первой краевой задачи для уравнения \hypertarget{ald:4*}{}

$$

L^*_1\upsilon \equiv \Delta_x \upsilon - \upsilon_t -

\sum\limits_{i=1}^m d_i \upsilon_{x_i}+ d\upsilon = \gamma

\upsilon

\eqno(4^*)

$$

с 

краевыми условиями

\begin{equation}

\upsilon \big|_{S} = \tau(r, \theta) = \bar{\tau}_n^k(r)  Y_{n,m}^k(\theta), \quad  \upsilon \big|_{\Gamma_\beta}

=0, 

 \label{ald:eq30}

\end{equation}

где  

$$d(x,t)= e -\sum\limits_{i=1}^m d_{ix_i},$$

$\bar{\tau}_{n}^k(r)\in G$, 

$G $ "--- множество функций $\tau(r)$ из класса $C\left(\left[0,1\right]\right) \cap C^1\left(\left(0,1\right)\right).$ 

Множество $G$ плотно всюду в $L_2 \left ((0,1)\right)$ \cite{aldash:9}.



Решение задачи (\hyperlink{ald:4*}{$4^\ast$}), \eqref{ald:eq30} будем искать в виде \eqref{ald:eq5}, где

функции $\bar{\upsilon}_n^k(r,t)$ будут определены ниже. 

Тогда, аналогично п.~\hyperlink{ald:p2}{2}, функции $\bar{\upsilon}_n^k(r,t)$ удовлетворяют 

системе уравнений \eqref{ald:eq7}--\eqref{ald:eq9}, где $\tilde{d}_{in}^k$, $d_{in}^k$

заменены на $-\tilde{d}_{in}^k$, $-d_{in}^k,$ а $\tilde{e}_n^k$ на $\tilde{d}_n^k,$ 

$i= 1,2, \dots,m$, $k =1,2, \dots, k_{n}$, $n=0,1,\dots.$



Далее в силу \eqref{ald:eq5} из краевых условий \eqref{ald:eq30}  получим

\begin{equation}

\bar{\upsilon}_n^k (r,0) = \overline{\tau}_n^k(r),\quad 

\bar{\upsilon}_n^k (1,t)=0,\quad k =1,2, \dots, k_{n},\quad n = 0,1,\dots .

 \label{ald:eq31}

\end{equation}



Как ранее замечено, каждое уравнение системы \eqref{ald:eq7}--\eqref{ald:eq9} представимо в виде \eqref{ald:eq10}. Задачу \eqref{ald:eq10}, \eqref{ald:eq31} сведем к следующей:

\begin{gather}

\label{ald:eq32}

L\upsilon_n^k \equiv  \upsilon_{nrr}^k -\upsilon_{nt}^k + \dfrac{\overline{\lambda}_n}{r^2}\upsilon_n^k - \gamma

\upsilon_n^k = f_n^k(r,t),

\\

\label{ald:eq33}

\upsilon_n^k(r,0)= \tau_n^k(r),\quad 

\upsilon_n^k(1,t) = 0,\quad k =1,2, \dots, k_{n},\quad n = 0,1,\dots, 

\\

\upsilon_n^k(r,t)= r^{\frac{(m-1)}{2}}\bar{\upsilon}_n^k (r,t),

\quad 

\tau_n^k(r)= r^{\frac{(m-1)}{2}}\bar{\tau}_n^k(r).

\end{gather}



Решение задачи \eqref{ald:eq32}, \eqref{ald:eq33} будем искать в виде

$$

\upsilon_{n}^k(r,t)=\upsilon_{1n}^k(r,t)+\upsilon_{2n}^k(r,t) ,

$$

где $\upsilon_{1n}^k(r,t)$ "--- решение задачи для \eqref{ald:eq32} с~условиями

\begin{equation}

\upsilon_{1n}^k(r,0)=0,\quad \upsilon_{1n}^k(1,t)=0, 

\label{ald:eq34}

\end{equation}

а $\upsilon_{2n}^k(r,t)$ "--- решение задачи для  уравнения

\begin{equation}

L \upsilon_{2n}^k = 0 

\label{ald:eq35}

\end{equation}

с условиями

\begin{equation}

\upsilon_{2n}^k(r,0)= \tau_n^k(r),\,\,\upsilon_{2n}^k(1,t)= 0.

\label{ald:eq36}

\end{equation}



Как показано в \cite{aldash:4}, решениями задач \eqref{ald:eq32}, \eqref{ald:eq34} и \eqref{ald:eq35}, \eqref{ald:eq36} являются, соответственно,  функции 

\begin{multline}

\upsilon_{1n}^k(r,t)= \sum\limits_{s=1}^\infty

\sqrt{r} \bigl(\exp(-\gamma - {\mu_{s,n}^2})t\bigr) \times \\

\times \biggl( \int_{t}^{0}a_{s,n}(\xi)(\exp(\gamma +

{\mu_{s,n}^2})d \xi\biggr)J_{n+{(m-2)}/{2}}(\mu_{s,n} r),

\end{multline}

$$

\upsilon_{2n}^k(r,t)= \sum\limits_{s=1}^\infty \tau_s

\sqrt{r}\bigl(\exp(-\gamma-{\mu_{s,n}^2})t\bigr)J_{n+{(m-2)}/{2}}(\mu_{s,n}

r),

$$

где

$$

\tau_s =

2[J_{n+{(m-2)}/{2}}(\mu_{s,n})]^{-2}\int_{0}^{1}\sqrt{\xi}\tau_n^k(\xi)J_{n+{(m-2)}/{2}}(\mu_{s,n}\xi)d\xi.

$$



Таким образом, решение задачи(\hyperlink{ald:4*}{$4^\ast$}), \eqref{ald:eq30}  построено в виде ряда

$$

\upsilon(r,\theta ,t) = {\sum\limits_{n = 0}^{\infty}

{{\sum\limits_{k =

1}^{k_{n}}r^{\frac{(1-m)}{2}}\bigl(\upsilon_{1n}^k(r,t)+ \upsilon_{2n}^k(r,t)\bigr)Y_{n,m}^{k} (\theta )}}}.

$$

При этом оно принадлежит классу $C(\overline{\Omega_\beta})\cap C^2(\Omega_\beta)$  \cite{aldash:4}.



Из определения сопряженных операторов $L_1$, $L^\ast_1$ \cite{aldash:11} имеем

$$

\upsilon L_1 u - u L^\ast_1 \upsilon = - \upsilon P(u) + u P(\upsilon) - u\upsilon Q,

$$

где

$$

P(u)=  \sum\limits_{i=1}^{m} u_{x_i} \cos \bigl(N^\perp, x_i\bigr), \quad  

Q = \cos\bigl(N^\perp, t \bigr)- \sum\limits_{i=1}^{m} d_{i} \cos \bigl(N^\perp, x_i\bigr),

$$

а $N^\perp$ "--- внутренняя нормаль к границе $\partial \Omega_\beta.$



Отсюда по формуле Грина получим

\begin{equation}

\int _S \tau(r,\theta)u (r,\theta,0)ds = 0. 

\label{ald:eq37}

\end{equation}



Поскольку линейная оболочка системы функций $\{\bar{\tau}_{n}^k(r)Y_{n,m}^k (\theta)\}$ плотна в $L_2(S)$~\cite{aldash:11}, из \eqref{ald:eq37} заключаем, что $u(r,\theta, 0)=0$ $~\forall

(r, \theta) \in S.$ Стало быть, по принципу экстремума для

уравнений \eqref{ald:eq4}~\cite{aldash:12} $u\equiv0$ в $\overline{\Omega}_\beta.$



В области $\Omega_\alpha$ получаем задачу \eqref{ald:eq28}, \eqref{ald:eq29},  для которой

имеет место теорема~\hyperlink{ald:th1}{1}~\cite{aldash:10}.

Теорема~\hyperlink{ald:th1}{1} доказана полностью.



\bigskip



\AdditionalContent{Конкурирующие интересы}{Конкурирующих интересов не имею.}



\AdditionalContent{Авторский вклад и ответственность}{Я несу полную ответственность за предоставление окончательной версии рукописи в печать. Окончательная версия 

рукописи мною одобрена.}





\AdditionalContent{Финансирование}{Исследование выполнялось без финансирования.}





\begin{thebibliography}{99}



\RBibitem{aldash:1}

\by  Нахушев~А.~М.

\book Задачи со смещением для уравнения в частных производных

\publaddr М.

\publ Наука

\yr 2006 

\totalpages 287 

\zmath{1135.35002}



\RBibitem{aldash:2}

\by  Врагов~В.~Н.

\book Краевые задачи для неклассических уравнений математической физики

\publaddr Новосибирск

\publ НГУ

\yr 1983

\totalpages 84



\RBibitem{aldash:3}

\by Алдашев~С.~А, Кожамкулов~Б.~А, Акитай~Б.~Е, Джумадиллаев~К.~Н.

\paper  Задачи Дирихле, возникающие при моделировании процессов деформации и разрушении композитов и их корректность \jour Вестник КазНПУ. Сер. Физ.-мат. науки

\yr  2014

\issue 2(46)

\pages 128--133



\RBibitem{aldash:4}

\by Алдашев~С.~А.

\paper  Корректность задачи Дирихле для одного класса многомерных гиперболо-параболических уравнений 

\jour Укр. матем. вестник

\yr 2013

\vol 10

\issue 2 

\pages 147--157



\RBibitem{aldash:5}

\by Алдашев~С.~А.

\paper Критерий однозначной разрешимости спектральной задачи Дирихле для многомерного

гиперболо-параболического   уравнения

\inbook Уравнения смешанного типа, родственные проблемы анализа и информатики

\bookinfo Второй Международный Российско-Узбекский симпозиум

\publaddr Нальчик

\publ НИИ ПМА КБНЦ РАН

\yr 2012

\pages 24--27



\RBibitem{aldash:6}

\by Михлин~С.~Г.

\book Многомерные сингулярные интегралы и интегральные уравнения

\publaddr М.

\publ Физматгиз

\yr 1962 

\totalpages 254 



\RBibitem{aldash:7}

\by Камке Э.

\book  Справочник по обыкновенным дифференциальным уравнениям

\publaddr М.

\publ Наука

\yr 1965 

\totalpages 703 



\Bibitem{aldash:8}

\by Erdélyi~A. (ed.)

\book  Higher transcendental functions

\bookvol II

\zmath{0505.33001}

\serial Bateman Manuscript Project, California Institute of Technology

\publaddr Malabar, Florida

\publ Robert E. Krieger Publ. 

\totalpages xviii+396 

\yr 1981



\RBibitem{aldash:9}

\by  Колмогоров~А.~Н., Фомин~С.~В.

\book  Элементы теории функций и функционального анализа

\publaddr М.

\publ Наука

\yr 1976 

\totalpages 543 



\RBibitem{aldash:10}

\by Алдашев~С.~А.

\paper Критерий однозначной разрешимости спектральной задачи Дирихле в~цилиндрической~области

для~многомерных гиперболических уравнений с~волновым оператором

\jour Вестн. Сам. гос. техн. ун-та. Сер. Физ.-мат. науки

\yr 2014

\issue 3(36)

\pages 21--30

\mathnet{http://mi.mathnet.ru/vsgtu1300}

\crossref{10.14498/vsgtu1300}

\elib{http://elibrary.ru/item.asp?id=23085708}



\RBibitem{aldash:11}

\by Смирнов~В.~И.

\book  Курс высшей математики

\bookvol 4, часть 2

\publaddr М.

\publ Наука

\yr 1981

\totalpages 550



\Bibitem{aldash:12}

\by  Friedman~A.

\book Partial differential equations of parabolic type

\zmath{0144.34903}

\publaddr Englewood Cliffs, N.J.

\publ Prentice-Hall, Inc. 

\totalpages xiv+347 

\yr 1964





\end{thebibliography}



\newpage



\makeentitle





\AdditionalContent{Competing interests}{I declare that I have no competing interests.} 

\AdditionalContent{Author's Responsibilities}{I take full responsibility for submitting the final manuscript in print. I approved the final version of the manuscript.} 

\AdditionalContent{Funding}{The research has not had any funding.}

%



\pagebreak





\begin{thebibliography}{99}



\Bibitem{aldash:1t}

\lang In Russian

\by  Nakhushev~A.~M.

\book Zadachi so smeshcheniem dlia uravneniia v chastnykh proizvodnykh \rm [Problems with Shifts for Partial Differential Equations]

\publaddr Moscow

\publ Nauka

\yr 2006 

\totalpages 287 



\Bibitem{aldash:2r}

\lang In Russian

\by  Vragov~V.~N.

\book Kraevye zadachi dlia neklassicheskikh uravnenii matematicheskoi fiziki \rm [Boundary-Value Problems for Nonclassical Equations of Mathematical Physics] 

\publaddr Novosibirsk

\publ  Novosibirsk State Univ.

\yr 1983

\totalpages 84



\Bibitem{aldash:3t}

\lang In Russian

\by Aldashev~S.~A, Kozhamkulov~B.~A, Akitai~B.~E, Dzhumadillaev~K.~N.

\paper  Dirichlet problems arising in the modeling of deformation and fracture processes of composites and their correctness 

\jour Vestnik KazNPU. Ser. Fiz.-mat. nauki

\yr  2014

\issue 2(46)

\pages 128--133



\Bibitem{aldash:4t}

\lang In Russian

\by Aldashev~S.~A.

\paper  Well-posedness of Dirichlet problem for one class of multi-dimensional hyperbolic-parabolic equations

\jour Ukr. Matem. Vestnik

\yr 2013

\vol 10

\issue 2 

\pages 147--157



\Bibitem{aldash:5t}

\lang In Russian

\by Aldashev~S.~A.

\paper A criterion for the unique solvability of the spectral Dirichlet problem for a multidimensional hyperbolic-parabolic equation

\inbook Equations of Mixed Type, Related Problems of Analysis and Informatics

\publaddr Nal'chik

\publ Research Institute of PMA KBSC RAS

\yr 2012

\pages 24--27



\Bibitem{aldash:6t}

\by Mikhlin~S.~G.

\book Multidimensional Singular Integrals and Integral Equations

\vol 83 

\serial International Series of Monographs on Pure and Applied Mathematics

\yr 1965

\zmath{0129.07701}

\publaddr Oxford, London, etc.

\publ Pergamon Press

\totalpages xii+255~pp \nofrills

\crossref{10.1016/C2013-0-01835-X}



\Bibitem{aldash:7t}

\lang In Russian

\by Kamke E.

\book  Spravochnik po obyknovennym differentsial'nym uravneniiam \rm [Handbook on Ordinary Differential Equations]

\publaddr Moscow

\publ Nauka

\yr 1965 

\totalpages 703 



\Bibitem{aldash:8t}

\by Erdélyi~A. (ed.)

\book  Higher transcendental functions

\bookvol II

\zmath{0505.33001}

\serial Bateman Manuscript Project, California Institute of Technology

\publaddr Malabar, Florida

\publ Robert E. Krieger Publ. 

\totalpages xviii+396 

\yr 1981



\Bibitem{aldash:9}

\lang In Russian

\by  Kolmogorov~A.~N., Fomin~S.~V.

\book  Elementy teorii funktsii i funktsional'nogo analiza \rm [Elements of the Theory of Functions and Functional Analysis]

\publaddr Moscow

\publ Nauka

\yr 1976 

\totalpages 543 



\Bibitem{aldash:10t}

\lang In Russian

\by Aldashev~S.~A.

\paper A criterion for the unique solvability of the Dirichlet spectral problem in a cylindrical domain for multidimensional hyperbolic equations with wave operator

\jour Vestn. Samar. Gos. Tekhn. Univ., Ser. Fiz.-Mat. Nauki \rm [J. Samara State Tech. Univ., Ser. Phys. Math. Sci.]

\yr 2014

\issue 3(36)

\pages 21--30

\crossref{10.14498/vsgtu1300}



\Bibitem{aldash:11t}

\by Smirnow~W.~I.

\book Course in higher mathematics

\bookinfo Part IV/2

\lang In German

\zmath{0997.00511}

\publaddr Frankfurt am Main

\publ H.~Deutsch

\totalpages 469 

\yr 1995



\Bibitem{aldash:12t}

\by  Friedman~A.

\book Partial differential equations of parabolic type

\zmath{0144.34903}

\publaddr Englewood Cliffs, N.J.

\publ Prentice-Hall, Inc. 

\totalpages xiv+347 

\yr 1964





\end{thebibliography}



\renewcommand{\baselinestretch}{0.9}



%

%\end{document}

